%LTeX: language=fr-FR
\documentclass[12pt,a4paper]{article}
\usepackage[utf8]{inputenc}
\usepackage[T1]{fontenc}
\usepackage[french]{babel}
\usepackage[margin=1.7cm]{geometry}
\usepackage{amsmath,amssymb}
\usepackage{graphicx}
\usepackage{enumitem}
\usepackage{tcolorbox}
\usepackage{xcolor}
\usepackage{multicol}
\usepackage{tasks}
\usepackage{tikz}
\usetikzlibrary{arrows.meta}

\tcbuselibrary{skins}

\definecolor{bandeau}{HTML}{E8EEF7}
\definecolor{titrebleu}{HTML}{1A4FA0}

% Bandeau de titre pour chaque exercice
\newtcbox{\exobox}{
  colback=bandeau, colframe=bandeau, boxrule=0pt,
  arc=0pt, left=4pt, right=4pt, top=3pt, bottom=3pt,
  fontupper=\bfseries, width=\linewidth, halign=left
}

% \exercice{numéro}{titre}{barème}
\newcommand{\exercice}[3]{%
  \vspace{0.15em}
  \noindent\exobox{Exercice #1 --- #2 \hfill (#3 points)}\par
  \vspace{0.15em}
}

\pagestyle{empty}
\setlist{topsep=2pt, itemsep=3pt, parsep=0pt}
\setlength{\parskip}{2pt}
\setlength{\parindent}{0pt}

% Barème détaillé indiqué entre crochets (version prof).
% La version élève (compile.sh, suffixe _ELEVE) retire le barème par question.
\ifdefined\ELEVE
\newcommand{\pts}[1]{}
\newcommand{\bareme}[1]{}
\else
\newcommand{\pts}[1]{\hfill\textcolor{titrebleu}{\small[#1]}}
\newcommand{\bareme}[1]{#1}
\fi

\colorlet{sol}{blue!70!black}
\colorlet{solJ}{red!70!black}

% Droite graduée de -10 à 10 ; #1 : tracé des intervalles
\newcommand{\droitecinq}[1]{%
  \begin{tikzpicture}[x=0.7cm,baseline=-0.5ex]
    \draw[{Latex[length=2mm]}-{Latex[length=2mm]}] (-10.6,0) -- (10.6,0);
    \foreach \i in {-10,...,10} {
      \draw (\i,0.08) -- (\i,-0.08);
      \node[below] at (\i,-0.08) {\scriptsize$\i$};
    }
    #1
  \end{tikzpicture}%
}
% \seg{couleur}{hauteur}{début}{fin}{crochet gauche}{crochet droit} (crochet vide : vers l'infini)
\newcommand{\seg}[6]{%
  \draw[#1, very thick] (#3,#2) -- (#4,#2);
  \node[#1, inner sep=0pt] at (#3,#2) {#5};
  \node[#1, inner sep=0pt] at (#4,#2) {#6};
}

\begin{document}

\begin{center}
  {\Large\bfseries Corrigé du rattrapage --- Intervalles}\\[0.4em]
  {\itshape Seconde \quad $\bullet$ \quad Contrôle de 50 minutes}
\end{center}

% ================= Exercice 1 =================
\exercice{1}{Automatismes}{5}
\begin{enumerate}
  \item On additionne les exposants : $3^{4}\times 3^{-2}=3^{4+(-2)}=3^{2}=\mathbf{9}$. \pts{1}
  \item $\dfrac{14}{15}-\dfrac{15}{16}<0$, donc $\dfrac{14}{15}<\dfrac{15}{16}$, c'est-à-dire $\mathbf{\dfrac{15}{16}>\dfrac{14}{15}}$. \pts{1}
  \item $P=6\times\left(3+\dfrac{2}{3}\right)=6\times 3+6\times\dfrac{2}{3}=18+4=\mathbf{22}$. \pts{1}
  \item On reconnaît $a^2+2ab+b^2$ avec $a=x$ et $b=4$ : $A=x^2+2\times x\times 4+4^2=\mathbf{(x+4)^2}$. \pts{2}
\end{enumerate}

% ================= Exercice 2 =================
\exercice{2}{Intervalles}{5}
\begin{enumerate}
  \item Appartenance. \bareme{$1$ point par réponse.} \pts{3}
  \begin{tasks}(2)
    \task $-2\notin\,\mathopen{]}-2\,;\,5]$ \quad ($-2$ est exclu)
    \task $6\in[1\,;\,6]$
    \task $\dfrac{3}{4}\notin[0{,}8\,;\,2]$ \quad ($\dfrac{3}{4}=0{,}75<0{,}8$)
  \end{tasks}
  \item Tableau. \bareme{$0{,}5$ point par case.} \pts{2}
  \begin{center}
  \renewcommand{\arraystretch}{1.4}
  \begin{tabular}{|c|c|c|c|c|}
  \hline
  Inégalité & $-1<x\leqslant 6$ & $x\leqslant -2$ & $\mathbf{-4\leqslant x<3}$ & $\mathbf{x>1{,}5}$ \\ \hline
  Intervalle & $x\in\mathbf{\mathopen{]}-1\,;\,6]}$ & $x\in\mathbf{\mathopen{]}-\infty\,;\,-2]}$ & $x\in[-4\,;\,3[$ & $x\in\,]1{,}5\,;\,+\infty[$ \\ \hline
  \end{tabular}
  \end{center}
\end{enumerate}

% ================= Exercice 3 =================
\exercice{3}{Intersection et union}{4}
$I$ en bleu, $J$ en rouge.

1. $I=[-3\,;\,+\infty[$ et $J=\,\mathopen{]}-6\,;\,2]$ : \quad $I\cap J=\mathbf{[-3\,;\,2]}$ \quad et \quad $I\cup J=\mathbf{\mathopen{]}-6\,;\,+\infty[}$. \pts{2}
\begin{center}
\droitecinq{\seg{sol}{0.3}{-3}{10.5}{$[$}{}\seg{solJ}{0.6}{-6}{2}{$]$}{$]$}}
\end{center}

2. $I=\,\mathopen{]}-\infty\,;\,-2]$ et $J=\,]1\,;\,5]$ : \quad $I\cap J=\mathbf{\varnothing}$ \quad et \quad $I\cup J=\mathbf{\mathopen{]}-\infty\,;\,-2]\,\cup\,]1\,;\,5]}$. \pts{2}

L'union n'est pas un intervalle.
\begin{center}
\droitecinq{\seg{sol}{0.3}{-10.5}{-2}{}{$]$}\seg{solJ}{0.6}{1}{5}{$]$}{$]$}}
\end{center}

% ================= Exercice 4 =================
\exercice{4}{Inégalités et inéquations}{6}
\begin{enumerate}
  \item \begin{enumerate}[label=\alph*)]
    \item $-3x+2\geqslant 14 \iff -3x\geqslant 12 \iff x\leqslant -4$ \quad (division par $-3<0$ : le sens change).

    $S=\mathbf{\mathopen{]}-\infty\,;\,-4]}$ \pts{2}
    \item $8-3x<x-4 \iff -4x<-12 \iff x>3$ \quad (division par $-4<0$ : le sens change).

    $S=\mathbf{]3\,;\,+\infty[}$ \pts{2}
  \end{enumerate}
  \item On multiplie par $-3<0$ : le sens des inégalités change. \pts{2}

  $x\in[-2\,;\,1[ \iff -2\leqslant x<1 \iff 6\geqslant -3x>-3 \iff 11\geqslant -3x+5>2$.\par Donc $x\in[-2\,;\,1[ \iff -3x+5\in\mathbf{]2\,;\,11]}$.
\end{enumerate}

\end{document}
